\writeSectionFrame{\presentationTitle}{Übung}

\begin{frame}
    \frametitle{Aufgabenbeschreibung}
    \begin{itemize}
        \item Implementieren Sie ein Pizza-Bestellsystem unter Verwendung des Dekorator-Musters.
        \item Die Basis-Pizza ist eine einfache Pizza ohne Beläge.
        \item Es gibt verschiedene Dekoratoren für die Pizza, wie z.B. Käse, Pepperoni und Oliven.
        \item Jeder Dekorator sollte die Beschreibung der Pizza erweitern und die Kosten entsprechend anpassen.
        \item Implementieren Sie eine abstrakte Klasse \texttt{Pizza}, die die Methoden \texttt{getDescription()} und \texttt{getCost()} definiert.
        \item Implementieren Sie eine konkrete Klasse \texttt{PlainPizza}, die eine einfache Pizza ohne Beläge repräsentiert.
        \item Erstellen Sie eine abstrakte Dekorator-Klasse \texttt{PizzaDecorator}, die von der \texttt{Pizza}-Klasse erbt und die Dekoratoren-Struktur umsetzt.
    \end{itemize}
\end{frame}

\begin{frame}
    \frametitle{Aufgabenbeschreibung}
    \begin{itemize}
        \item Erstellen Sie mindestens drei konkrete Dekoratoren-Klassen, z.B. \texttt{CheeseDecorator}, \texttt{PepperoniDecorator}, und \texttt{OlivesDecorator}.
        \item Testen Sie das System in der \texttt{main}-Methode, indem Sie verschiedene Pizzen mit unterschiedlichen Belägen erstellen und die Beschreibung und Kosten ausgeben.
    \end{itemize}
\end{frame}

\begin{frame}
    \frametitle{Erweiterungen (Optional)}
    \begin{itemize}
        \item Fügen Sie eine Kalorienberechnung für jeden Belag hinzu.
        \item Jeder Dekorator sollte die Kalorien der Basis-Pizza entsprechend erweitern.
        \item Stellen Sie sicher, dass die Kalorien korrekt für jede Pizza-Bestellung berechnet werden.
    \end{itemize}
\end{frame}